\section{The description entropy of all binary qubit effects}\label{sec:biasedentropy75}
Let
\[
 \mathfrak E=\{E=E^*\in\C^{2\times2}:0\le E\le I\}.
\]
An effect $E$ specifies the ordered measurement with effects $E,I-E$.
Write $\mathcal C_N(\mathfrak E,\delta)$ for its covering number under
$d_N$, allowing arbitrary legal memoryless centres of this interface,
and write $\mathcal C_N^{\mathrm{rat}}(\mathfrak E,\delta)$ when the
centres have rational Cartesian coordinates. Both eigenvalues and the
eigenspaces are charged. The spectral representation
$E=pP_u+q(I-P_u)$, $p\ge q$, represents a scalar effect only once,
independently of $u$.

\begin{theorem}[The full effect body]\label{thm:biasedcover75}
There are absolute constants $c,C,\delta_0>0$ such that, for every integer
$N\ge1$ and $0<\delta\le\delta_0$,
\begin{equation}\label{eq:biasedcover75}
 cN^2\log(N+2)\,\delta^{-4}
 \le \mathcal C_N(\mathfrak E,\delta)
 \le \mathcal C_N^{\mathrm{rat}}(\mathfrak E,\delta)
 \le CN^2\log(N+2)\,\delta^{-4}.
\end{equation}
Consequently the optimum fixed-length reusable payload is
\begin{equation}\label{eq:biasedbits75}
 2\log_2N+\log_2\log(N+2)+4\log_2(1/\delta)+O(1).
\end{equation}
The constants are uniform over the whole closed effect body, including
scalar, deterministic and projective effects. For rational accuracy
input, rational target data admit an exact rational encoder attaining
this order.
\end{theorem}
The logarithm records accumulation near the projective stratum, where
both outcome probabilities can approach their support endpoints.
We prove the upper bound by a finite rational construction and the
lower bound by separated families of projective-corner boxes.

\subsection{A rational spectral and angular code}
For $N\ge1$ and rational $0<\delta\le1/4$, put
\begin{equation}\label{eq:biasedspectralgrid75}
 B=\left\lceil\sqrt{64N/\delta^2}\right\rceil,\qquad
 s_i=\begin{cases}
 \displaystyle\frac{i^2}{B^2+i^2},&0\le i\le B,\\[2mm]
 \displaystyle1-\frac{(2B-i)^2}{B^2+(2B-i)^2},&B<i\le2B.
 \end{cases}
\end{equation}
Thus $0=s_0<s_1<\cdots<s_{2B}=1$. For $0\le j<i\le2B$ define
\begin{equation}\label{eq:biasedangulargrid75}
 K_{ij}^2=
 \frac{N(s_i-s_j)^2}
 {s_i(1-s_i)+s_j(1-s_j)+N^{-1}},\qquad
 A_{ij}=\max\left\{1,
      \left\lceil\sqrt{256K_{ij}^2/\delta^2}\right\rceil\right\}.
\end{equation}
Use the six signed stereographic charts
\begin{equation}\label{eq:biasedsphere75}
 u_a=\epsilon\frac{1-|z|^2}{1+|z|^2},\qquad
 u_b=\frac{2z_b}{1+|z|^2}\quad(b\ne a),\qquad
 a\in\{1,2,3\},\quad \epsilon\in\{-1,1\},
\end{equation}
with $z\in\{-A_{ij},\ldots,A_{ij}\}^2/A_{ij}$ for the spectral
pair $(s_i,s_j)$. Every chart word decodes to
$s_iP_u+s_j(I-P_u)$. A diagonal pair $i=j$ has the single word
$s_iI$. The exact number of words, including redundant chart words,
is
\begin{equation}\label{eq:biasedcapacity75}
 L_{N,\delta}=(2B+1)+
       \sum_{0\le j<i\le2B}6(2A_{ij}+1)^2.
\end{equation}

\begin{theorem}[Exact rational realization]\label{thm:biasedcodec75}
The words in \eqref{eq:biasedcapacity75} form a radius-$\delta/2$
cover of $\mathfrak E$ by legal rational effects, and
\begin{equation}\label{eq:biasedcapacitybound75}
 L_{N,\delta}\le C N^2\log(N+2)\,\delta^{-4}
\end{equation}
for an absolute $C$. If
\[
 E=\tfrac12\bigl((1+b)I+x\cdot\sigma\bigr),\qquad
 b\in\Q,\quad x\in\Q^3,\quad |b|+|x|\le1,
\]
then a covering word is computable by rational arithmetic and integer
comparisons. Its payload is one fixed-length index in the set of
\eqref{eq:biasedcapacity75}.
\end{theorem}
\begin{proof}
Write $f(t)=t^2/(1+t^2)$, $0\le t\le1$. For a probability in
$[0,1/2]$, round its inverse coordinate $t$ to the nearest multiple
of $1/B$; for a probability in $[1/2,1]$, use the reflected chart
$1-f(t)$. Choose the multiple of smaller magnitude at ties. This
defines a nondecreasing quantizer on $[0,1]$, so separately quantizing
$p\ge q$ preserves their order.

The square-root vector associated with $f(t)$ is
$(t,1)/\sqrt{1+t^2}$; its derivative has norm $1/(1+t^2)\le1$.
The reflected chart has the same bound. Hence the Euclidean
square-root distance between each probability and its quantization
is at most $1/(2B)$. If $H$ is this distance, the product-fidelity
inequality gives unhalved distance at most $2\sqrt N H$ between the
corresponding $N$-fold Bernoulli laws. The common two-coin programme
in Lemma~\ref{lem:spectralprojection75} therefore gives a spectral error
at most
\[
                         2\sqrt N/B\le\delta/4.
\]
This argument includes adaptive testers and entangled references.

If the quantized eigenvalues coincide, this completes the encoding.
Otherwise choose a signed chart whose distinguished coordinate has
maximal absolute value in the target direction. Its inverse coordinate
belongs to $[-1,1]^2$. Rounding its two coordinates to the grid of
denominator $A_{ij}$ gives angular error at most $2\sqrt2/A_{ij}$,
since geodesic distance is bounded by the length of the mapped
straight segment and the differential of stereographic projection has
norm at most two. The angular upper bound in Theorem~\ref{thm:biasedangular75}
gives error at most
\[
 \sqrt2 K_{ij}\frac{2\sqrt2}{A_{ij}}
                   \le\delta/4.
\]
The triangle inequality proves the covering assertion. Formula
\eqref{eq:biasedsphere75} preserves the unit sphere exactly, so every
decoded effect is legal and rational.

We verify exact computability even when the input eigenvalues are
irrational. Put $Q=|x|^2\in\Q$. The eigenvalues are
$(1+b\pm\sqrt Q)/2$. Each decision in the spectral rounding compares
one of these numbers with a rational threshold $f((k+1/2)/B)$ or its
reflection. Such a comparison is the sign of a rational number plus
or minus $\sqrt Q$; checking signs before squaring makes it an exact
rational comparison. If $Q=0$, the two quantized eigenvalues coincide.
Otherwise choose the smallest index $a$ maximizing $|x_a|$ and put
$\epsilon=\operatorname{sgn}(x_a)$. The inverse chart is
\[
                    z_b=\frac{x_b}{\sqrt Q+|x_a|}\quad(b\ne a).
\]
Its comparison with any rational rounding threshold again reduces
to the sign of a rational number plus a rational multiple of
$\sqrt Q$. The integers $B,A_{ij}$ are computed by integer squaring
and rational comparison. Thus no irrational normalization is required.
The decoder uses only the displayed rational formulae.

It remains to count the words. Set $m_i=\min(i,2B-i)$ and
$U=B^2/N$. Directly from \eqref{eq:biasedspectralgrid75},
\[
 \frac{m_i^2}{4B^2}\le s_i(1-s_i)\le\frac{m_i^2}{B^2},\qquad
 K_{ij}^2\le
          \frac{4NB^2}{m_i^2+m_j^2+U}.
\]
Each value of $m_i$ occurs at most twice. Since
$(2A_{ij}+1)^2\le C(1+K_{ij}^2/\delta^2)$, we obtain
\begin{equation}\label{eq:biasedlattice75}
 L_{N,\delta}\le CB^2+
 \frac{CNB^2}{\delta^2}
       \sum_{a,b=0}^{B}\frac1{a^2+b^2+U}.
\end{equation}
Here $U\ge1$. The part with $\max(a,b)\le\sqrt U$ is bounded by an
absolute constant. Each successive dyadic square ring contributes
at most another absolute constant. There are at most
$C\log(N+2)$ such rings, because $B/\sqrt U=\sqrt N$. Consequently
\[
              \sum_{a,b=0}^{B}\frac1{a^2+b^2+U}
                       \le C\log(N+2).
\]
Finally $B^2\le CN/\delta^2$, so \eqref{eq:biasedlattice75} proves
\eqref{eq:biasedcapacitybound75}. The support cutoff in the denominator
is retained throughout the count, including the endpoint grid words.
\end{proof}

\subsection{Packing the projective corner}
\begin{proof}[Proof of Theorem~\ref{thm:biasedcover75}]
The upper bound follows from Theorem~\ref{thm:biasedcodec75}. For a
real $\delta$, choose a rational $\delta'$ with
$\delta\le\delta'\le2\delta$ and use its radius-$\delta'/2$ code;
decrease $\delta_0$ to at most $1/8$. For the lower bound first let
$N\ge16$. For every
\[
                     t=8^j/N\le1/16,\qquad j\ge0,
\]
consider the effects with
\begin{equation}\label{eq:biasedcorner75}
 p=1-\alpha,\quad q=\beta,\quad
 t\le\alpha,\beta\le2t,\qquad
 u=u(z),\quad z\in[-1/4,1/4]^2,
\end{equation}
where $u(z)$ is the chart in \eqref{eq:biasedsphere75} with
$a=3,\epsilon=1$. On this square, spherical distance is bounded
above and below by absolute positive multiples of Euclidean distance
in $z$. Throughout \eqref{eq:biasedcorner75}, the eigenvalue gap is
at least $3/4$ and
\[
 p(1-p)\asymp q(1-q)\asymp t,\qquad
 p(1-p)+q(1-q)+N^{-1}\asymp t.
\]
The finite-distance comparison of Theorem~\ref{thm:biasedmetric75}
therefore supplies an absolute $c_1>0$ such that any two points in
this box satisfy
\begin{equation}\label{eq:biasedboxmetric75}
 d_N\ge c_1\min\left\{1,
       \sqrt{N/t}\,
       \bigl| (\alpha,\beta,z)-
              (\alpha',\beta',z')\bigr|_\infty\right\}.
\end{equation}
Choose an absolute sufficiently large $A$ and then an absolute
sufficiently small $\delta_0$. In each of the four coordinates take
a grid of spacing $A\delta\sqrt{t/N}$. Since $Nt\ge1$, this spacing
is at most half the width of either spectral interval when
$\delta\le\delta_0$. The resulting points have pairwise distance
at least $4\delta$ by \eqref{eq:biasedboxmetric75}, and their number
is at least
\[
 c\left(\frac{t}{\delta\sqrt{t/N}}\right)^2
   \left(\frac1{\delta\sqrt{t/N}}\right)^2
                         =cN^2\delta^{-4}.
\]

Points from different boxes are also uniformly separated. Indeed,
if $t'\ge8t$, then $\alpha'\ge t'$ and $\alpha\le2t\le t'/4$.
The spectral modulus of the corresponding upper eigenvalues is
bounded below by an absolute positive multiple of
$\sqrt{Nt'}\ge1$. The spectral noncancellation bound in
Theorem~\ref{thm:biasedmetric75} gives $d_N\ge c_2>0$.
Shrinking $\delta_0$ ensures that this is at least $4\delta$.
There are at least $c\log(N+2)$ permitted boxes. Their union is
therefore a $4\delta$-packing of cardinality
$cN^2\log(N+2)\delta^{-4}$.

For $1\le N<16$, use instead a fixed spectral rectangle
$5/8\le p\le3/4$, $1/4\le q\le3/8$ and the same angular chart.
Theorem~\ref{thm:biasedmetric75} gives a uniform local Euclidean lower
bound, and a grid of spacing $C\delta$ has at least $c\delta^{-4}$
points. The factor $N^2\log(N+2)$ is bounded on this finite range.
This proves the same lower estimate for all $N\ge1$.

A radius-$\delta$ ball about any legal centre contains at most one
point of a $4\delta$-packing, by the triangle inequality. Hence the
packing proves the lower bound with the stated centre convention.
Taking binary logarithms and rounding the index length proves
\eqref{eq:biasedbits75}.
\end{proof}

\subsection{The source of the logarithm}
The spectral and angular weights also give the following integral
form of the same accumulation. Put $v_p=p(1-p)$, $v_q=q(1-q)$ and
$\tau=N^{-1}$. Then
\begin{equation}\label{eq:biasedvolume75}
 \int_{0\le q\le p\le1}
 \frac{N^2(p-q)^2\,dp\,dq}
 {(v_p+v_q+\tau)\sqrt{(v_p+\tau)(v_q+\tau)}}
                         \asymp N^2\log(N+2).
\end{equation}
To verify this, near $p=1,q=0$ put $\alpha=1-p$, $\beta=q$.
After removing the factor $N^2$, the integrand is comparable to
\[
 \frac1{(\alpha+\beta+\tau)
             \sqrt{(\alpha+\tau)(\beta+\tau)}}.
\]
For $s=\alpha+\beta$, integration over $0\le\alpha\le s$ is
bounded above by an absolute constant times $1/(s+\tau)$.
For $s\ge2\tau$, restriction to $s/4\le\alpha\le3s/4$ gives the
matching lower bound. Integration in $s$ proves logarithmic growth.
Near the scalar endpoint $p=q=0$, the factor $(p-q)^2$ is at most
$(p+q)^2$, while $v_p+v_q\asymp p+q$; the remaining integral is
uniformly bounded. Complementation gives the same conclusion near
$p=q=1$. Away from these endpoint neighbourhoods,
$v_p+v_q$ is bounded below and the factors $v_p^{-1/2},v_q^{-1/2}$
are integrable. Bounded $N$ is absorbed into the constants.
Thus the extra logarithm is contributed by the projective corner.
The finite covering construction and packing above also include the
collapsed scalar coordinates, where the spectral representation
ceases to be a chart.
